(1) 每个入射光子能量为 $h\nu$，最少必须提供金属表面的逸出功 $\Phi$ 才能释放电子，因此最大发射动能满足Einstein方程 $K_{\max}=h\nu-\Phi$，仅在 $h\nu\ge\Phi$ 时有实际发射。多余能量转为电子平移动能，而不是决定电子发射总数。
(2) 截止条件是 $h\nu_c=\Phi$，所以 $\nu_c=\Phi/h$。对应真空波长 $\lambda_c=c/\nu_c=hc/\Phi$；对于 $\lambda>\lambda_c$ 的光，即使增加单色弱场照明时间，单个光子仍不足以克服逸出功，理想单光子光电发射被禁止。
(3) 遏止电压定义为反向电势提供的能量 $eV_s=K_{\max}$，故 $V_s=(h/e)\nu-\Phi/e$（$\nu>\nu_c$）。斜率 $dV_s/d\nu=h/e$ 与金属逸出功无关，截距则由 $\Phi$ 决定，这可用于实验测量Planck常数。
(4) 每秒入射光子数为 $P/(h\nu)$，若独立光子导致电子计数的量子效率为 $\eta$，则光电子平均发射率 $\dot N_e=\eta P/(h\nu)$，光电流 $I=e\dot N_e=e\eta P/(h\nu)$。此式要求光子能量超过阈值；低于阈值的理想单光子效率应为零。
